\documentclass{jsarticle}
\usepackage[dvipdfmx,final]{graphicx}
\usepackage{chemfig}
\usepackage{amsmath}
\usepackage{amssymb}
\begin{document}
\title{Time scale measure replaced with space scale volume of \\
global deprivation and integral manifold \\
scale volume is space ideality measure}
\author{Masaaki Yamaguchi}
\date{}
\maketitle

\newcommand{\variint}{%
	\begin{picture}(15,30)
		\put(0,0){
			$ \iint $}
		\put(0,5){\line(15,0){15}}
	\end{picture} %
}
   \newcommand{\squareiint}{%
	\begin{picture}(15,30)
		\put(0,0){
			$ \iint $}
		\put(0,5){$\text{\scalebox{2.0}[1]{$\square$}}$}
	\end{picture} %
}

  
   \newcommand{\dazanier}{%
	\begin{picture}(15,30)
		\put(0,0){
			$ \bigsqcup $}
		\put(4,3){$\text{\scalebox{1.0}[1]{$\top$}}$}
		\put(4,3){$\text{\scalebox{1.0}[1]{$\bot$}}$}
	\end{picture} %
}


\newcommand{\variiint}{%
	\begin{picture}(15,15)
		\put(0,0){
			$ \iiint $}
		\put(0,5){\line(15,0){20}}
	\end{picture} %
}



 \newcommand{\alearletter}{%
	 \begin{picture}(10,20)
		 \put(0,0){
			 $ \square $}
		 \put(0,0){\line(1,1){10}}
	 \end{picture} %
	 }

  \newcommand{\whitewithbodercircle}
  {\text{\textcircled{\phantom{\textbullet}}}
  \kern-0.9em\text{\textcircled{\phantom{}}}}
     

 
  \newcommand{\whitewithblack}{\text{\textcircled{\textbullet}}}
  \newcommand{\whitewithblacksquare}{\text{\fbox{\text{\rule{0.5em}
  {0.5em}}}}}
  \newcommand{\whitewithwhitesquare}{\text{\fbox{\text{\phantom
  {\rule{0.5em}{0.5em}}}}}}
 \newcommand{\whitewithbodercircleiiint}{%
	 \begin{picture}(30,45)
		 \put(0,0){\scalebox{2.0}[2]{$\whitewithbodercircle$}}
		 \put(0,0){\scalebox{2.0}[2]{$\iiint$}}
	 \end{picture}%
	 }
\newcommand{\whitewithbodercircleint}{%
	 \begin{picture}(30,45)
		 \put(0,0){\scalebox{2.0}[2]{$\whitewithbodercircle$}}
		 \put(0,0){\scalebox{2.0}[2]{$\int$}}
	 \end{picture}%
	 }
 \newcommand{\whitewithbodercirclevarint}{%
	 \begin{picture}(30,45)
		 \put(0,0){\scalebox{2.0}[2]{$\whitewithbodercircle$}}
		 \put(0,0){\scalebox{2.0}[2]{$\varint$}}
	 \end{picture}%
	 }
 
 \newcommand{\whitewithbodercirclevariiint}{%
	 \begin{picture}(30,45)
		 \put(0,0){\scalebox{2.0}[2]{$\whitewithbodercircle$}}
		 \put(0,0){\scalebox{2.0}[2]{$\variiint$}}
	 \end{picture}%
	 }
 
 \newcommand{\average}{%
	  \begin{picture}(30,45)
		  \put(0,0){\scalebox{2.0}[2]{$\square$}}
		  \put(4,4){$\setminus$}
	  \end{picture}%
  } 
 
 These equation are Projection function, cross time, assemble of projection
 function, estimate manifold, global integral manifold, project global 
 integral manifold, average of add and even function.
 This system is signal of write with system.
  $$ \whitewithblack f_m , f \boxtimes g = F^{m}_t, 
  \whitewithblacksquare f_m = {\begin{pmatrix}
	  n \\
  r \end{pmatrix}}^{\otimes L} = \lim_{n \to \infty}
  {\begin{pmatrix}
	  n \\
  r \end{pmatrix}}^{\otimes L} = {d \over {d f}}\boxtimes $$

  \newcommand{\whitewithwhitecircle}
  {\text{\textcircled{\phantom{\textbullet}}}}
  Estimate function is equal with dervargence function
  $$ \begin{pmatrix}
	  n \\
  r \end{pmatrix}(\zeta(s)) = \whitewithbodercircle f_m $$
  $$ \whitewithblack f_m ,
  \whitewithblacksquare f_m = {\begin{pmatrix}
	  n \\
  r \end{pmatrix}}^{\otimes L} = \lim_{n \to \infty}
  {\begin{pmatrix}
	  n \\
  r \end{pmatrix}}^{\otimes L} = {d \over {d f}}\boxtimes $$


 

$$ \whitewithbodercircleint \sum [dI_m] $$
$$ {}^t\variiint_{D\chi} cohom D\chi = \log(x\log x) = [x_m,y_m]\times
[x_n,y_n] $$
$$ \int R^{\nabla r}dr_x = l^{R|\nabla} = \beta(p,q) $$
$$ V_{D\chi}\int dx_m = S \otimes S = \vee\int dx_m = S^2 \otimes S^2 $$
$$ \Delta[\vartriangleleft_x,\vartriangleleft_y] = \average(x,y) $$

  \newcommand{\whitewithboder}{%
	  \begin{picture}(15,30)
		  \put(0,0){\scalebox{2.0}[2]{$\square$}}
		  \put(4,4){$\square$}
	  \end{picture}%
  } 
  \newcommand{\whitewithnabla}{%
	  \begin{picture}(15,30)
		  \put(0,0){\scalebox{2.0}[2]{$\nabla$}}
		  \put(3,3){$\nabla$}
	  \end{picture}%
  } 


  Assemble function is Euler equation, and this equation also
  Euler product of fundermantal manifold.
  $$ \whitewithnabla = e^{i\theta} = \cos \theta + i \sin \theta $$

  Average function.
  $$ \average = {{f(x) + g(x)}} \ge 2{\sqrt{f(x)g(x)}} $$
  Global manifold is partial integral of global topology.
 $$ \whitewithblacksquare f_m =  \whitewithboder \text{                 
      } f_m = 
 \whitewithbodercircle f_m = 
  D_{\mathcal{K}}(x) = \int dx_m , \mathrm{ker}/\mathrm{im}f = 
  \partial $$

$$\boxtimes$$
$$\whitewithbodercircle, \whitewithblack $$
$$\whitewithblacksquare,\whitewithwhitecircle$$
$$\scalebox{2.0}[2]{\variint}$$ 
$$\whitewithboder$$ 
$$ \whitewithnabla$$
$$\average$$ 
は、それぞれ、積のクロス、 大域的積分多様体、組み合わせ多様体、
同型の大域的積分多様体、準同型の大域的積分多様体、
準同型写像、大域的切断、重ね合わせの原理、
相加相乗平均、をそれぞれ表す。
 

$$||ds^2||=\whitewithbodercircleiiint \mathrm{Docol}[I_m]$$
$$  \whitewithbodercircleint \sum[dI_m]$$

$$ {}^t\scalebox{2.0}[2]{\variint}_{D\chi}\mathrm{colm}D\chi = \log(x\log x)$$
$$ = [x_m,y_m]\times [x_n,y_n] $$

\newcommand{\timebow}{%
	\begin{picture}(5,5)
		\put(0,0){$\text{\rotatebox{90}{$\bowtie$}}$}
	\end{picture} %
	}


 \newcommand{\flowerletter}{%
	 \begin{picture}(10,20)
		 \put(0,0){$\text{\scalebox{1.0}[1]{$\square$}}$}
		 \put(0,0){$\text{\scalebox{1.0}[2]{$\int$}}$}
	 \end{picture} %
	 }




     ヒッグス場方程式が、相加相乗平均方程式であり、大域的多様体でもある。
  $$ {d \over {df}}F(x,y)=m(x) $$
  $$ = {d \over {d m(x)}}M = \partial (\sigma(M)) $$
  $$=  {d \over {d \timebow}}Z^{'}(\zeta)dx_m $$
  $$ ||ds^2|| = \nabla_{i}\nabla_{j}\int \nabla M(\sigma)d\eta $$
  M理論の指数定理による擬微分と擬積分多様体が、ヘルマンダー作用素から
  のノルム空間となり、
  $$ = M^{{m(x)^{'}}} $$
  M理論のニュートン方式の大域的微分多様体になる。
  $$ ||ds^2|| = \int_M [d(\partial) + d(\sigma)]dI_m $$
  それが、D-braneのM理論へと導かれる。
  $$ {d \over {df}}(e^{-f} + e^{f}) = {{d \over {d R^{'}}}}R = d[I_m] $$
  結局は、Jones多項式における、大域的多様体としての、虚数微分方程式となる。
  リッチ・フロー方程式による、大域的多様体の擬微分と擬積分多様体の
  計算式が、ニュートン方式による、擬積分と擬微分多様体となり、
  時間計量が、相対性による、測量がなくなり、空間計量となり、
  空間概念の量子化となる。
$$ ||ds^2|| = \whitewithbodercircleiiint Docol[I_m] $$

 	
$$ \timebow = \flowerletter r dr_m $$

 
$$ \whitewithbodercircleint \sum [dI_m] $$
$$ {}^t\variiint_{D\chi} cohom D\chi = \log(x\log x) = [x_m,y_m]\times
[x_n,y_n] $$
$$ \int R^{\nabla r}dr_x = l^{R|\nabla} = \beta(p,q) $$
$$ V_{D\chi}\int dx_m = S \otimes S = \vee\int dx_m = S^2 \otimes S^2 $$
$$ \Delta[\vartriangleleft_x,\vartriangleleft_y] = \average(x,y) $$


         時間計量による、指数定理における擬微分と擬積分多様体を微小変量によって、
   大域的トポロジーが成されられる。ガンマ関数における大域的多様体と、
   ヒッグス方程式における大域的多様体の計算式を以下に、載せておく。
   $$ {T^{\mu\nu}}^{T^{\mu\nu}} = \int T^{\mu\nu} \rotatebox{90}{$\bowtie$}_m $$
   $$   {T^{\mu\nu}}^{{T^{\mu\nu}}^{'}}=\int T^{\mu\nu}dx_m  $$
      $$ {\Gamma^{\gamma}}^{'} = {d \over {d \gamma}}\Gamma $$
   $$ \Gamma^{\gamma} = \int \Gamma dx_m $$
   $$ \rotatebox{90}{$\bowtie$} = \int x^x dx_m $$


$$ ||ds^2|| = \whitewithbodercircleiiint Docol[I_m] $$

 	
$$ \timebow = \flowerletter r dr_m $$


  始まりは、無限の時間の無が、場の理論として存在していた。
  この無が、オイラーの定数の大域的積分多様体が、
  基本群による、相対性から、オイラーの定数の多様体積分としての
  ガンマ関数に化けて、ゼウスとなり、このゼウスから、
  質量差エネルギーとして、母なるJones多項式ができた。
  このJones多項式から、各種の生成物として、
  宇宙と異次元が、D-braneを柱として、生み出された。
  このJones多項式から、中間子を媒介にして、素粒子が出来て、
  世界が切り開かれた。
  この世界は、Zeta関数を禅に対して、Beta関数が終わりに来るように、
  時間が２通りの流れとして、存在された。
  この時間の流れの中で、互いの相対性が消えて、
  空間概念の量子化としての、すべての共時性が最終目標にされた。


     Zata function escourt into Forgotfull of underlying modify on Sum
   of summative group, thiis Forgotfull summative group instented with
   Space ideal theorem from quantum level of deprivate space in aspect
   of quantum level of differential geometry,
   This theorem construct with Higgs field and Morse theory, moreover
   Hortshorn conjecture, fourth step deduction of theorem from 
   Euler product of integral manifold.
   This manfold revolte with global differential manifold in global topology.
   From these theorem inspirate from Forgotfull theorem in Space ideality
   theory, and this theory conclude into entropy of non exchange equation.
   Forgotten theory estimate with zeta function oneselves.
   Forgotten theory income with non relativity theory and this theory
   face to face without partner and partner of non relativity from
   ideal of space in revolution theory.
   All exist and non exist partner with non relativity escourt into 
   Forgotten theory, and this theory include with same underly group
   of all of partner with same distance ideal.

   $$ X \dotsi \to Y, Y \dotsi \to X $$
   $$ \bigoplus(i \hbar^{\nabla})^{\oplus L} = e^{-x \log x} $$
   This equation of quantum level of differential geomerty
   instimate with entorpy of non exchange equation.
   This equation means with low level of botton entropy in Space ideal of
   Forgotten theory. Moveover, this equation enterstein with all of exist
   theory from general relativity.
   No time and No Space of relativity, and monotonicity of magnetic component
   with gravity and antigravity, and this theory comment with partner of
   magnetic component theory included with being resulted from Forgotten
   theory.
   $$ \square \to \nabla |: x \to y, f(x) \to g(y), f^{-1}(x)xf(x) 
   \dotsi \to g(y) $$
   $$ \square \dotsi \to \nabla |: f(x) \cong g(y) $$
   $$ {d \over dt}g_{ij}(t) = - 2 R_{ij}, F^m_t = -2\int{{(R + \nabla_i \nabla
   _j f)} \over {(\Delta + f)}}dm $$
   $$ {d \over df}F = m(x) $$
   $$ F|_{\_} = {d \over df}F, F^m_t = e^{-f}dV $$
   $$ \bigoplus(i \hbar^{\nabla})^{\oplus L} = e^{-x \log x} $$
   These equation represent with Forgotten theory.
   空間概念の量子化は、代数幾何の量子化であり、忘却関数でもあり、
   忘却関手のカテゴリー論でもあり、簡潔に言うと、すべての相対性が
   機能すると、全て同型というコホモロジー論でのコイコライザーである。
   このコイコライザーが、空間概念の量子化である。
   すべては、ゼータ関数へと行き着く。
   ベータ関数は、誤差関数であり、宇宙の雑音として存在していると、
   異次元へと移行する。
   $$ \beta(p,q) \cong {\Gamma(p)\Gamma(q) \le \Gamma(p+q)}, 
   {{\Gamma(p)\Gamma(q)} \over \Gamma(p+q)} \le 1 $$
   $$ \int e^{-x}x^{1-t}dx_m = \int e^{-x}x^{1-t}dx\cdot e^{-\theta} $$
   $$ = \Gamma^{{\gamma}^{'}} $$
   $$ = e^{-x \log x} $$
   閉３次元多様体と３次元球面が、空間概念の量子化としての結果であり、
   この多様体が、無としてのベータ関数を形作っている３次元多様体である。
   宇宙と異次元合わせてが、空間概念の量子化での単体である。
   コイコライザーとは、共変微分によって対になっている変数に作用する写像であり、
   このコイコライザーで処理された変数同士を同型にさせる群が、
   同型のコイコライザーである。そして、最小の単体でもある。
   故に、宇宙と異次元合わせて、コイコライザーである。
   球対称である一般相対性理論の平坦な空間が、空間概念の量子化でもある。
   ビッグ・バンは、イコライザーであり、
   $$ \langle f,g\rangle \mapsto \langle d,d\rangle $$
   として、宇宙と異次元の曲平面を表している。


$$ \dazanier = D_{\chi}{}^{t}\scalebox{1.0}[2]{\variiint}
\mathrm{exp}{[{\mathrm{cohom}}D_{\chi}]}[I_m] $$ 

　Scopion sigunature is globale cutten assemble of tuple space,
and this space is belong with Cohomology value.
$$ D_{\chi} |: x \to y $$
D-brane of variable is projection of pole of real and imaginary extranention.
$$ D_{\chi}(\chi)y = \log D_{\chi}^{y} = x $$
Regular of matrix of group theory is set theory included.
$$ \square \cap \Psi \ni  \zeta(s) $$
$$ \int \kappa (A^{\mu\nu})^{2}dx_m = \int \Gamma^{'}(\gamma)dx_m $$
Gravity theory equal in gamma of global partial integral manifold.
$$ g_{\mu\nu}(x)g^{-1}_{\mu\nu} = {}^{t}\scalebox{1.0}[2]{\variiint}_{D\chi}
\kappa (A^{\mu\nu})dx_m $$

Kalut-klein theory involved with regure group value with einsieten equation.
$$  \mathcal{S}^{\mu\nu}_{D(\chi,D)} = \bigoplus {(V_{k}{}^{N}\squareiint
\Gamma(\gamma)d\gamma)}[I_m] $$
Sheap theory have on own equation of metamorphose of global connected manifold.
$$ \mathcal{O}(x) = \oint \pi x^2 dx $$
$$ \int e^{-x}x^{1-t}dx = \int \sin \theta \cos \theta d \theta \int e^{\sin \theta \cos \theta}d\theta $$

 These equation are resulted from zeta function reconsturct with Higgs and
 Beta function metamorphose of Gut theory.

\end{document}
